The fresh L21 studies identify a prediction-error component that is highly oracle-correctable in the tested setting, then demonstrate learned usefulness on separate fresh data. The unchanged historical rank-eight space supports a 99.25\% oracle readout-MSE reduction; a learned coefficient correction supports a 97.21\% reduction relative to the historical predictor on another 64 blocks. These are different cohorts and interventions on prediction error, not an equality between oracle and learned performance. The space was fitted from native vector responses without direct readout labels. Separate rank/norm controls favor FIT-derived spaces over the fixed random draws; ranks 16 and 32 improve further, so neither an intrinsic dimension eight nor optimality is established.

The baseline comparison is decisive for interpretation. Direct384 has lower mean loss than learned correction despite using the same additional 384 vector labels, and a linear direct predictor restricted to $U$ also has slightly lower point loss. A constant mean correction worsens the historical predictor. The fresh budget TEST additionally identifies a conditional C32 advantage of 48.03\% over Direct32, but Direct has lower mean loss at 64--384 labels (the 256/384 contrast intervals include zero) and C64 worsens relative to C32. Three new FIT32 cohorts on the reused TEST favor Direct on average for have/has; the is/are mean interval includes zero. The C32 advantage belongs to one fixed training sequence; two of the later is/are fits have negative point contrasts, but the mixed three-fit result does not establish stable superiority. Thus the evidence supports scoped learned utility and an observed conditional advantage. C is prior-centered kernel ridge, with an exact linear prior interpretation under the conditions in Section~\ref{sec:priorridge}; it is not a new or inherently more expressive regression algorithm. Historical training information is not equalized by matching the additional label count. Native prefix computation is still needed.

The L27/L17 and earlier tests establish limits rather than contradicting this L21 result. The same three Qwen L27 fitted pairs help on T2 and harm on T5; Gemma-270M benefits on both, while separate Gemma-1B and Llama-1B cohorts establish T5 harm; Llama benefits on T2. Their operators, layers and fitting recipes differ, preventing a causal attribution to layer alone. In L27, accurate oracle sign forecasts coexist with a gate that worsens Qwen T5, and Direct256 has lower mean loss on both frames. Repairability, learned utility and advantage over direct prediction therefore remain distinct empirical questions.

Most conclusions concern one is/are margin in synthetic contexts; the posthoc have/has extension tests another grammatical margin on the same prompts, selected checkpoints/operators and fixed fits. Fresh pairs do not undo historical selection of the favorable L21 regime. Nominal intervals do not establish fit-population stability or familywise claims; fixed random controls do not sample the entire distribution of possible spaces. The empirical contribution is the controlled sequence from oracle-correctable error to measured learned utility and its limits, with negative baseline comparisons retained. Public release still requires complete authorship, evidence availability and campaign disclosures.
